\section{Related Work}
\label{sec:related}

Classical organization-centered multi-agent systems define roles
independently of their occupants, through formal specifications such as
Moise$^+$ and OperA and through norms that institutions regiment or
sanction~\cite{hubner2002moise,dignum2004opera,esteva2001formal,grossi2007ubilex}.
An LLM role prompt, in contrast, is a degenerate case: it states such
constraints in prose but has no mechanism to enforce them. Prompted roles
are accordingly known to be unreliable~\cite{cemri2025why,orogat2026architectural}.
Moreover, an evolved Planner, Executor and Critic team is no better than an
evolved single agent when the number of model calls is held
fixed~\cite{dylan2026equal}, and many reported multi-agent efficiency gains
arise where one agent suffices~\cite{zhang2026rethinking}.

Iterative refinement, tool feedback and multi-agent debate are standard
paradigms~\cite{madaan2023selfrefine,gou2024critic,du2024debate}, yet models
often repeat earlier outputs even under exact verifier
feedback~\cite{jiang2026exactfeedback}, and self-correction degrades
accuracy when the model judges its own answer~\cite{huang2024selfcorrect}.
It does help when an oracle's correct label
serves as a stopping rule~\cite{huang2024selfcorrect}, and our loop is
gated by exactly such an oracle, an evaluator that knows the correct
answer, where more attempts solve more problems~\cite{brown2024monkeys}. The closest study separates reviewer
precision from critique uptake~\cite{yang2026precise} but neither
manipulates the role structure nor equalizes feedback access across
conditions; we do both (Appendix~\ref{app:related}).
